\documentclass[../DoAn.tex]{subfiles}
\begin{document}

\label{chapter:2}

\section{Tổng quan}
Chương 1 đã trình bày bối cảnh, bài toán và định hướng xây dựng hệ thống JobBridge AI. Để hiện thực hóa một nền tảng tuyển dụng có tích hợp trí tuệ nhân tạo, đồ án cần dựa trên nhiều nhóm công nghệ khác nhau, từ kiến trúc phần mềm, phát triển giao diện, xây dựng API, lưu trữ dữ liệu, đến triển khai và vận hành trên môi trường đám mây. Chương này trình bày các cơ sở lý thuyết và công nghệ nền tảng được sử dụng trong hệ thống, đồng thời làm rõ vai trò của từng lựa chọn công nghệ đối với bài toán tuyển dụng thông minh.

Nội dung của chương được tổ chức theo các nhóm vấn đề chính: kiến trúc microservices cho backend, công nghệ phát triển ứng dụng phía server và client, nền tảng container hóa và điều phối bằng Kubernetes, quy trình CI/CD theo mô hình GitOps, và cuối cùng là cơ sở lý thuyết về mô hình ngôn ngữ lớn được ứng dụng trong các chức năng AI của hệ thống.

\section{Kiến trúc Microservices}
\label{section:2.1}
Microservices là phong cách kiến trúc trong đó hệ thống được chia thành nhiều dịch vụ nhỏ, độc lập, mỗi dịch vụ đảm nhiệm một miền nghiệp vụ rõ ràng và giao tiếp với nhau thông qua API. Khác với kiến trúc nguyên khối, nơi toàn bộ chức năng được đóng gói trong một ứng dụng duy nhất, microservices cho phép từng thành phần được phát triển, kiểm thử, triển khai và mở rộng một cách tương đối độc lập. Cách tiếp cận này phù hợp với JobBridge AI vì hệ thống có nhiều nhóm chức năng khác nhau như xác thực người dùng, quản lý tin tuyển dụng, quản lý hồ sơ ứng tuyển và xử lý AI.

Trong hệ thống JobBridge AI, phần backend được tổ chức thành bốn dịch vụ chính viết bằng Go. API Gateway chạy ở cổng 8080 đóng vai trò điểm vào duy nhất của frontend và thực hiện định tuyến request theo đường dẫn đến các service phía sau. Auth Service chạy ở cổng 8081 xử lý đăng ký, đăng nhập, phát hành và xác thực JWT, quản lý hồ sơ người dùng, công ty, avatar và CV. Jobs Service chạy ở cổng 8082 xử lý danh sách việc làm, tin tuyển dụng của nhà tuyển dụng và hồ sơ ứng tuyển. AI Service chạy ở cổng 8085 phụ trách các chức năng Interview Coach, Interview Quiz và đánh giá CV bằng mô hình ngôn ngữ lớn. Cả Auth, Jobs và AI Service đều truy cập MongoDB để lấy dữ liệu nghiệp vụ cần thiết; riêng AI Service còn gọi OpenAI-compatible API để sinh nội dung AI.

\begin{figure}[H]
    \centering
    \makebox[\textwidth][c]{%
        \includegraphics[width=1.15\textwidth]{Chuong/so-do-chuong-2/microservices_architecture.png}%
    }
    \caption{Sơ đồ kiến trúc microservices của JobBridge AI}
    \label{fig:microservices-architecture}
\end{figure}

Điểm quan trọng trong thiết kế hiện tại là API Gateway không trực tiếp xác thực JWT hay quyết định quyền nghiệp vụ. Gateway chủ yếu đóng vai trò reverse proxy, còn việc xác thực và phân quyền được thực hiện tại từng service thông qua middleware và logic xử lý nghiệp vụ. Ví dụ, các chức năng AI dành cho ứng viên yêu cầu người dùng có vai trò ứng viên và JWT hợp lệ; chức năng đánh giá CV yêu cầu người dùng có vai trò nhà tuyển dụng, đồng thời kiểm tra nhà tuyển dụng có sở hữu tin tuyển dụng liên quan hay không. Cách tổ chức này giúp mỗi service tự bảo vệ tài nguyên của mình và hạn chế phụ thuộc vào giả định rằng mọi request đều đi qua một lớp gateway duy nhất.

Kiến trúc microservices mang lại ba lợi ích chính cho đồ án. Thứ nhất, hệ thống được tách theo miền nghiệp vụ nên dễ bảo trì và dễ mở rộng. Thứ hai, mỗi service có thể được đóng gói thành Docker image riêng và triển khai độc lập trên Kubernetes. Thứ ba, khi một nhóm chức năng có tải cao hơn, chẳng hạn Jobs Service trong quá trình tìm kiếm việc làm hoặc AI Service trong quá trình gọi mô hình, hệ thống có thể mở rộng riêng service đó mà không cần nhân bản toàn bộ ứng dụng.

\section{Các công nghệ sử dụng}
\label{section:2.2}
Việc lựa chọn công nghệ trong JobBridge AI dựa trên ba tiêu chí: phù hợp với yêu cầu nghiệp vụ tuyển dụng, hỗ trợ phát triển nhanh trong phạm vi đồ án tốt nghiệp, và đủ khả năng mở rộng khi triển khai trên môi trường container hóa. Hình \ref{fig:tech-stack} minh họa tổng quan các công nghệ chính được sử dụng trong hệ thống.

\begin{figure}[H]
    \centering
    \makebox[\textwidth][c]{%
        \includegraphics[width=1.15\textwidth]{Chuong/so-do-chuong-2/4_technology_stack.png}%
    }
    \caption{Các công nghệ chính được sử dụng trong dự án}
    \label{fig:tech-stack}
\end{figure}

\subsection{Công nghệ phía Backend}
\label{subsection:2.2.1}
Backend của hệ thống được phát triển bằng Go, một ngôn ngữ biên dịch có hiệu năng tốt, cú pháp gọn và khả năng đóng gói thuận lợi cho môi trường container. Go phù hợp với các dịch vụ REST API vì có thư viện chuẩn mạnh cho lập trình mạng, đồng thời hỗ trợ xử lý đồng thời bằng goroutine. Trong đồ án, Go được sử dụng thống nhất cho Gateway, Auth Service, Jobs Service và AI Service nhằm giảm độ phức tạp khi phát triển, kiểm thử và triển khai.

Gin Gonic được sử dụng làm web framework cho các service backend. Framework này cung cấp cơ chế routing, middleware, binding dữ liệu đầu vào và trả response JSON một cách đơn giản. Trong hệ thống, Gin được dùng để tổ chức các nhóm API theo miền nghiệp vụ như xác thực, hồ sơ người dùng, việc làm, hồ sơ ứng tuyển và chức năng AI. Các middleware xác thực JWT và kiểm tra vai trò được áp dụng ở những route cần bảo vệ, trong khi các route công khai như xem danh sách việc làm hoặc xem công ty có thể truy cập mà không cần đăng nhập.

MongoDB được chọn làm cơ sở dữ liệu chính vì dữ liệu tuyển dụng có cấu trúc linh hoạt, bao gồm hồ sơ người dùng, thông tin công ty, tin tuyển dụng, đường dẫn CV, nội dung CV đã trích xuất và trạng thái ứng tuyển. Hệ thống sử dụng các collection chính gồm người dùng, công ty, tin tuyển dụng và hồ sơ ứng tuyển. Các quan hệ như người dùng sở hữu công ty, nhà tuyển dụng sở hữu tin tuyển dụng hoặc ứng viên ứng tuyển vào một vị trí được quản lý ở tầng ứng dụng thông qua mã định danh logic, thay vì dùng ràng buộc khóa ngoại như trong cơ sở dữ liệu quan hệ.

Đối với lưu trữ tệp, hệ thống sử dụng Cloudinary để lưu avatar và CV. Khi ứng viên tải CV lên, Auth Service gửi tệp lên Cloudinary, nhận về đường dẫn truy cập, đồng thời trích xuất nội dung văn bản từ PDF và lưu lại trong hồ sơ người dùng. Cơ chế này giúp AI Service có thể lấy nội dung CV trực tiếp từ MongoDB trong các lần xử lý sau, giảm chi phí tải lại và phân tích lại tệp.

\subsection{Công nghệ phía Frontend}
\label{subsection:2.2.2}
Frontend được xây dựng bằng React kết hợp TypeScript và Vite. React hỗ trợ tổ chức giao diện thành các component có thể tái sử dụng, phù hợp với hệ thống có nhiều trang theo vai trò như public user, ứng viên, nhà tuyển dụng và admin. TypeScript giúp kiểm soát kiểu dữ liệu ở tầng giao diện và API client, giảm lỗi trong quá trình phát triển. Vite được sử dụng làm công cụ build và dev server, giúp rút ngắn thời gian khởi động và hỗ trợ hot module replacement trong quá trình lập trình.

Về giao diện, hệ thống sử dụng Tailwind CSS để xây dựng bố cục và định dạng theo hướng nhất quán. Các biểu tượng được cung cấp bởi thư viện Lucide React. Ứng dụng frontend được tổ chức theo các nhóm chức năng như xác thực, việc làm, công ty, khu vực ứng viên, khu vực nhà tuyển dụng và khu vực quản trị. Cách tổ chức này giúp tách biệt logic gọi API theo miền nghiệp vụ và giao diện hiển thị theo vai trò người dùng.

\section{Nền tảng triển khai và DevOps}
\label{section:2.3}
Hệ thống JobBridge AI được thiết kế không chỉ như một ứng dụng web đơn lẻ mà còn như một hệ thống có thể triển khai và vận hành trên môi trường Kubernetes. Vì vậy, các công nghệ DevOps đóng vai trò quan trọng trong việc bảo đảm tính nhất quán môi trường, khả năng tái tạo và khả năng cập nhật tự động.

\subsection{Container hóa với Docker}
\label{subsection:2.3.1}
Docker được sử dụng để đóng gói từng thành phần của hệ thống thành container image. Mỗi service backend và frontend đều có cấu hình đóng gói riêng, giúp ứng dụng chạy nhất quán giữa môi trường phát triển, CI và Kubernetes. Với backend Go, image có thể được tối ưu bằng multi-stage build: giai đoạn đầu biên dịch chương trình, giai đoạn sau chỉ chứa tệp thực thi và các thành phần runtime tối thiểu. Với frontend, giai đoạn build tạo ra các tệp tĩnh và Nginx phục vụ kết quả đó ở môi trường production.

\subsection{Điều phối với Kubernetes}
\label{subsection:2.3.2}
Kubernetes là nền tảng điều phối container được sử dụng để chạy các service của JobBridge AI trên Azure Kubernetes Service. Trong Kubernetes, mỗi service được triển khai dưới dạng Deployment và được expose nội bộ bằng Service. Frontend và Gateway được truy cập từ bên ngoài thông qua Ingress, trong khi Auth, Jobs, AI và MongoDB chủ yếu phục vụ traffic nội bộ trong cluster. Helm chart được dùng để quản lý các manifest Kubernetes theo dạng template, cho phép tái sử dụng cấu hình giữa môi trường local và Azure.

\begin{figure}[H]
    \centering
    \makebox[\textwidth][c]{%
        \includegraphics[width=1.22\textwidth]{Chuong/so-do-chuong-2/kubernetes_deploymeny.png}%
    }
    \caption{Kiến trúc triển khai trên Azure Kubernetes Service}
    \label{fig:kubernetes-deployment}
\end{figure}

Trong triển khai Azure, hệ thống sử dụng Azure Container Registry để lưu trữ image, Azure Key Vault để quản lý bí mật như JWT secret, OpenAI API key và Cloudinary URL, đồng thời sử dụng Secrets Store CSI Driver để đưa các bí mật này vào pod. Các tài nguyên Kubernetes quan trọng gồm Deployment, Service, Ingress, StatefulSet cho MongoDB, Horizontal Pod Autoscaler, Pod Disruption Budget và SecretProviderClass.

\subsection{CI/CD với GitHub Actions và Argo CD}
\label{subsection:2.3.3}
Quy trình CI/CD của hệ thống được thiết kế theo hướng GitOps. Khi mã nguồn thay đổi, GitHub Actions thực hiện các bước kiểm tra chất lượng, chạy kiểm thử, build image, quét bảo mật image bằng Trivy và đẩy image lên Azure Container Registry. Image được gắn tag theo mã commit để có thể truy vết phiên bản triển khai. Sau đó, pipeline triển khai cập nhật cấu hình Helm dùng cho Argo CD để trỏ đến phiên bản image mới.

\begin{figure}[H]
    \centering
    \makebox[\textwidth][c]{%
        \includegraphics[width=1.15\textwidth]{Chuong/so-do-chuong-2/ci_flow.png}%
    }
    \caption{Sơ đồ luồng tích hợp liên tục của dự án}
    \label{fig:ci-flow}
\end{figure}

\begin{figure}[H]
    \centering
    \makebox[\textwidth][c]{%
        \includegraphics[width=1.15\textwidth]{Chuong/so-do-chuong-2/cd_flow.png}%
    }
    \caption{Sơ đồ luồng triển khai liên tục của dự án}
    \label{fig:cd-flow}
\end{figure}

Argo CD chạy trong AKS và liên tục so sánh trạng thái mong muốn trong Git với trạng thái thực tế của cluster. Khi phát hiện cấu hình Helm thay đổi, Argo CD đồng bộ tài nguyên Kubernetes và triển khai phiên bản mới theo cơ chế rolling update. Mô hình GitOps giúp toàn bộ thay đổi triển khai được ghi nhận trong lịch sử Git, dễ kiểm tra, dễ rollback và hạn chế thao tác thủ công trực tiếp trên cluster.

\section{Trí tuệ nhân tạo và Mô hình ngôn ngữ lớn}
\label{section:2.4}
Trí tuệ nhân tạo là lĩnh vực nghiên cứu và xây dựng các hệ thống có khả năng thực hiện những tác vụ thường đòi hỏi trí thông minh của con người. Trong phạm vi đồ án này, trọng tâm không phải là huấn luyện mô hình AI từ đầu mà là ứng dụng mô hình ngôn ngữ lớn thông qua API để giải quyết các bài toán xử lý ngôn ngữ tự nhiên trong tuyển dụng.

Mô hình ngôn ngữ lớn, hay Large Language Model, là mô hình được huấn luyện trên khối lượng dữ liệu văn bản rất lớn, có khả năng hiểu ngữ cảnh, sinh văn bản, tóm tắt, so sánh và trả lời theo chỉ dẫn. Với JobBridge AI, LLM được sử dụng trong ba tình huống chính: hỗ trợ ứng viên luyện phỏng vấn theo JD và CV, sinh bộ câu hỏi trắc nghiệm phỏng vấn theo vị trí đã ứng tuyển, và hỗ trợ nhà tuyển dụng đánh giá mức độ phù hợp của CV ứng viên so với mô tả công việc.

Để kết quả AI bám sát dữ liệu thực tế, hệ thống xây dựng prompt từ nhiều lớp ngữ cảnh. Đối với ứng viên, prompt bao gồm hồ sơ cá nhân, mô tả công việc, yêu cầu công việc, nội dung CV đã trích xuất và lịch sử hội thoại gần nhất. Đối với nhà tuyển dụng, prompt bao gồm thông tin ứng viên, JD, CV text và yêu cầu đánh giá. Bên cạnh đó, hệ thống áp dụng các kiểm tra nghiệp vụ trước khi gọi mô hình, chẳng hạn chỉ ứng viên đã ứng tuyển mới được sử dụng AI Coach cho job đó, và chỉ nhà tuyển dụng sở hữu job mới được đánh giá CV của ứng viên trong job tương ứng. Điều này giúp hạn chế việc sử dụng AI sai ngữ cảnh và tăng tính an toàn dữ liệu.

\section{Kết chương}
Chương này đã trình bày các nền tảng lý thuyết và công nghệ chính được sử dụng trong JobBridge AI. Kiến trúc microservices giúp hệ thống tách biệt các miền nghiệp vụ và mở rộng linh hoạt; Go, Gin, React, TypeScript và MongoDB là các công nghệ cốt lõi để xây dựng ứng dụng; Docker, Kubernetes, Helm, GitHub Actions, Argo CD, Terraform và Azure Key Vault tạo nên nền tảng triển khai và vận hành; mô hình ngôn ngữ lớn cung cấp khả năng phân tích và sinh phản hồi thông minh cho các chức năng tuyển dụng. Các nội dung này là cơ sở để chuyển sang Chương 3, nơi hệ thống được phân tích và thiết kế chi tiết hơn theo các yêu cầu nghiệp vụ.

\end{document}
